\begin{figure}[htbp]
\centering
\setlength{\unitlength}{1mm}
\begin{picture}(150,65)
  \put(2,45){\makebox(38,6)[l]{\small View of the plane $z=0$}}
  \put(3,7){\framebox(66,37){}}
  \put(26,16){\framebox(20,20){\large $D$}}
  \put(28,10){\makebox(16,5){\scriptsize fixed square}}
  \put(6,37){\makebox(19,5)[l]{\scriptsize $\Pi\setminus D$}}
  \put(4,1){\makebox(65,5)[l]{\scriptsize plane continues beyond frame}}
  \put(78,45){\makebox(70,6)[l]{\small Power accounting}}
  \put(79,37){\vector(1,0){12}}
  \put(76,41){\makebox(21,4){\scriptsize input $a$}}
  \put(92,31){\framebox(26,14){\shortstack{\scriptsize dielectric\\[-1pt]\scriptsize device}}}
  \put(119,37){\vector(1,0){12}}
  \put(131,33){\makebox(18,8)[l]{\shortstack[l]{\scriptsize selected\\[-1pt]\scriptsize output $Ta$}}}
  \put(79,19){\makebox(68,7)[l]{\scriptsize On $D$: signed forward flux}}
  \put(79,11){\makebox(68,7)[l]{\scriptsize Outside $D$: absolute flux}}
  \put(79,3){\makebox(68,7)[l]{\scriptsize Inequalities hold for every $a\in W$}}
\end{picture}
\caption{Schematic of the channel definition. The left panel is a view of
the whole measurement plane, with a fixed square $D$; the outer frame only
marks the drawing boundary. The right panel indicates input--output power
accounting, not a prescribed one-way device geometry. Passage uses signed
flux through $D$, whereas bypass measures the absolute flux on the rest of
the plane. In the counterexample the selected output is the complete
outgoing scattering map, including reflection.}
\label{fig:channel-geometry}
\end{figure}
